\subsubsection{A: Extension-Bending-Torsion}
From \eqref{eq:section-linear-tangent}, it is evident that the complete response of a linear isotropic prism is characterized by the constant geometric cross section tensors:
\begin{equation*}
\left[
\begin{matrix}
    A \mathbf{1} 
  & A\bar{\bm{\SectionLocation}}^{\times} % nm
  & \SectionNW  % E
  & \SectionNV  % G
  \\

%
  & \LinearMM
  & \LinearMW % E
  & \LinearMV % G
  \\

%
  & %
  & \LinearWW % E?
  & \\

%
  & %
  & %
  & \LinearVV
\end{matrix}
\right]
\triangleq 
\SectionIntegral{
\TotalStrainMatrix^{\top} \TotalStrainMatrix
}
= \SectionIntegral{
\left[
\begin{matrix}
\mathbf{1} 
  &-\bm{r}^{\times} 
  & \FrameBasis\otimes \bm{\varphi} 
  & \FrameBasis_{\beta} \otimes \nabla_{\beta} \bm{\varphi} \\

%- \bm{r}^{\times} 
  & - \left[\bm{r}^{\times}\right]^{2} 
  & \bm{r} \times \FrameBasis \otimes \bm{\varphi} 
  & \bm{r} \times \FrameBasis_{\beta} \otimes \nabla_{\beta} \bm{\varphi} \\

%\bm{\varphi} \otimes \FrameBasis
  & %[\bm{\varphi} \otimes \FrameBasis] \bm{r}^{\times} 
  & \bm{\varphi} \otimes \bm{\varphi}
  & \mathbf{0} \\

%
  & % 
  & %
  & \nabla_{\beta} \bm{\varphi} \otimes \nabla_{\beta} \bm{\varphi}
\end{matrix}
\right]
}
% \label{eq:def-linear-isotropic-coupling-tensors}
\end{equation*}
where \(A\) is the cross sectional area, \(A \bar{\bm{\SectionLocation}}\) are the first moments of area, \(\LinearMM\) is the standard tensor of second area moments, \(\LinearWW\) generalizes the warping constant of \cite{vlasov1961thinwalled} to \(n_w\) warping shapes and \(\LinearVV\) encodes Saint Venant's constant as well as the shear correction factors of \cite{timoshenko1921correction}, when the warping shapes \(\bm{\varphi}\) are chosen appropriately.

\begin{Quote}
$$
\begin{aligned}
\FrameBasis_{\beta}\otimes\nabla_{\!\FrameBasis_{\beta}}\bm{\psi}
&=\left[\begin{array}{cc}
 \psi_{y,y} & \psi_{z,y} \\
 \psi_{y,z} & \psi_{z,z} \\
\end{array}\right]_{\FrameBasis_{\alpha}\otimes\FrameBasis_{\beta}}
\end{aligned}
$$

$$
\TotalStrainMatrix_{\rm p} = \left[\begin{array}{ccc|ccc}
1 & 0 & 0 &  0 & z & -y \\
0 & 1 & 0 & -z & 0 &  0 \\
0 & 0 & 1 &  y & 0 &  0 \\
\end{array}\right]
$$
\end{Quote}


\paragraph{$\bm{n}$-$\bm{n}$}

\paragraph{m-m}\label{m-m}

\switchocg{tensor index}{
\begin{ocg}{jmm-tensor}{tensor}{on}
\[ 
\boldsymbol{\sigma}=\mathbf{C}:\boldsymbol{\varepsilon}
\]
\end{ocg}
\begin{ocg}{jmm-index}{index}{off}
\[
  \begin{bmatrix}
    \sigma_{xx}\\[2pt] \sigma_{yy}\\[2pt] \sigma_{xy}
  \end{bmatrix}
  =
  \begin{bmatrix}
    \lambda+2\mu & \lambda & 0 \\
    \lambda & \lambda+2\mu & 0 \\
    0 & 0 & \mu
  \end{bmatrix}
  \begin{bmatrix}
    \varepsilon_{xx}\\[2pt] \varepsilon_{yy}\\[2pt] 2\varepsilon_{xy}
  \end{bmatrix}
\]
\end{ocg}
}

\[
\LinearMM = \SectionIntegral{
     [\bm{r}^\times][\bm{r}^\times]^{\top} 
}
= \SectionIntegral{\bm{r}\cdot\bm{r}\mathbf{1} - \bm{r}\otimes\bm{r}}
% &= \int (\bm{r}\cdot\bm{r})\mathbf{1} - \bm{r}\otimes \bm{r}\, \mathrm{d}A \\
% &=\left[\delta_{\alpha \beta} \mathbf{1}-\FrameBasis_\alpha \otimes 
\]

\paragraph{w-w}\label{w-w}

\begin{equation}
\LinearWW = \SectionIntegral{\bm{\varphi}\otimes\bm{\varphi}}
= \SectionIntegral{ \begin{bmatrix}
\varphi^2  & \varphi \psi_y & \varphi \psi_z \\
           & \psi_y^2 & \psi_z \psi_y \\
           &          & \psi_z^2
\end{bmatrix}}
= \begin{bmatrix}
\Jw & 0 & 0 \\
0 & C_{\psi} & 0 \\
0 & 0 & C_{\psi}
\end{bmatrix}
\end{equation}

\paragraph{\(\bm{v}-\bm{v}\)}

\begin{equation}
\LinearVV = \SectionIntegral{\nabla_{\beta} \bm{\varphi}\otimes\nabla_{\beta} \bm{\varphi}}
\end{equation}
\textcite{gruttmann2000theory} refer to \(\Jv\) 
% NOTE: I previously used C_{\alpha} for J_v
with the notation ``\(I_0 + I_T\)''
\[
\Jv  = \SectionIntegral{ \WarpTwistSC_{, y}^2 + \WarpTwistSC_{, z}^2} = J_S - J
\]


\begin{proposition}
(Twist centre). The twist centre $\mathrm{C}^{\rm tw}$ is the pole in the plane of the cross section $\Omega$, such that a tangential stress field $\bm{\tau}\left(\bm{b}_{\Section}, c_{\vartheta}\right)$ with vanishing resultant moment about $\mathrm{C}^{\rm tw}$ performs no mutual work when interacting with any twist tangential strain field $\bm{\varepsilon}_{\gamma}^{\rm tw}(c_{\vartheta}, \bm{r})=c_{\vartheta}\left(\FrameBasis\times\bm{r}+\nabla \phi^{\rm tw}\right)$. 
The position of $\mathrm{C}^{\rm tw}$ is given by the formula:

$$
\begin{aligned}
& \bm{r}_{\mathrm{C}^{\rm tw}}=\FrameBasis \times\,\IntRGoRG^{-1} \int_{\Omega} \phi^{\rm tw}(\CoordinateCentroid ) \CoordinateCentroid \mathrm{~d}A 
\qquad\text{ where }\quad 
\phi^{\rm tw}:=\varphi^{\rm tw}(1) \in \mathrm{C}^2(\Omega ; \mathcal{R})
\end{aligned}
$$
\end{proposition}


\paragraph{Coupling in $\bm{n}-\bm{m}$}\label{n-m}
Static moment:
\[
\LinearNM
\triangleq \SectionIntegral{ -[\bm{r}^\times] } \equiv -A {\CentroidLocation}^{\times}
\]
where \textbf{two} independent constants known as the \emph{first moments of area} have been defined.

\paragraph{Coupling in $\bm{n}-\bm{w}$}\label{n-w}
Sectorial static moment:
\[
\SectionNW \triangleq \SectionIntegral{ \FrameBasis \otimes\bm{\varphi} }
\] 
where \textbf{three} independent constants have been defined
\paragraph{Coupling in n-v}
\[
\SectionNV \triangleq \SectionIntegral{ \FrameBasis_{\beta}\otimes \nabla_{\beta}\bm{\varphi}
}
= \begin{bmatrix}
0 &     0 & 0 \\
J_{ny} & S_{nyy} & S_{nzy} \\
J_{nz} & S_{nyz} & S_{nzz}
\end{bmatrix}
\] 
assuming \(\bm{\varphi}' = \mathbf{0}\).

\paragraph{m-w}\label{m-w}
% NOTE: May remove the symbols J_my, etc... if they are always equal to -Iv, then they are redundant.
Sometimes referred to as the \emph{Sectorial deviation moment}:
\[
\LinearMW 
= \SectionIntegral{ \bm{r} \times\FrameBasis \otimes \bm{\varphi} } 
= \begin{bmatrix}
0 & 0 & 0 \\
J_{my} & S_{myy} & S_{mzy} \\
J_{mz} & S_{myz} & S_{mzz}
\end{bmatrix}
\]
introducing \textbf{two} shear-interacting constants \(J_{my}\) and \(J_{mz}\) for torsion.
%
For moment-bishear interaction:
\[
\LinearMV
= \SectionIntegral{ \bm{r} \times\FrameBasis_{\beta} \otimes \nabla_{\beta}\bm{\varphi} }
\]
Generally, however, \(\LinearMV \equiv - \LinearVV\).
\[
\LinearMW 
= \SectionIntegral{ \bm{r} \times\FrameBasis \otimes \bm{\varphi} } 
= \SectionIntegral{ \begin{bmatrix}
0 & -z \varphi_{, y} + y \varphi_{, z} \\
 z \varphi & 0 \\
-y \varphi & 0
\end{bmatrix}
}
= \begin{bmatrix}
0 & S_{\alpha} \\
S_{y} & 0 \\
S_{z} & 0
\end{bmatrix}
\]



\begin{remark}\label{rem:torsion-constant}
The torsion constant \(\Jsv\) integrates the warping function \(\WarpTwistIesan\) which considers twist about the origin of the coordinate system.

The classical Saint-Venant torsion constant \(\Jsv\) integrates the warping function \(\WarpTwistSC\) defined about the shear center. 
For \(\Jsv\), the following identities hold:
\begin{equation*}
    \begin{aligned}
        \Jsv &= \int (\FrameBasis\times\tilde{\bm{r}} + \nabla \WarpTwistSC) \cdot (\FrameBasis\times\tilde{\bm{r}} + \nabla \WarpTwistSC) \mathrm{~d}A \\
        &= \int \FrameBasis\times\tilde{\bm{r}} \cdot (\FrameBasis\times\tilde{\bm{r}} + \nabla \WarpTwistSC) \mathrm{~d}A \\
        &= \int \tilde{\bm{r}}\cdot\tilde{\bm{r}} + \FrameBasis\times\tilde{\bm{r}} \cdot \nabla \WarpTwistSC \\
        &= \left<\tilde{\bm{r}},\, \tilde{\bm{r}}\right>_G + \left\{\WarpTwistSC,\, \tfrac12 r^2\right\}_G
    \end{aligned}
\end{equation*}
where the following was employed:
\[
\int \nabla \WarpTwistSC \cdot (\FrameBasis\times \tilde{\PlanePosition} + \nabla\WarpTwistSC)\mathrm{~d}A = 0.
\]
Define the constants \(\Jv\) and \(\Jw\):
\begin{equation*}
    \Jv = ( \WarpTwistSC,\,\WarpTwistSC)_G,
    \qquad
    J_{\rm{mv}} = \left\{\WarpTwistSC,\,\tfrac12 r^2\right\}_G
    \qquad\text{and}\qquad 
    \Jw = \left<\WarpTwistSC,\,\WarpTwistSC\right>_E
\end{equation*}
Then the following identities hold:
\begin{equation*}
\begin{aligned}
    \Jsv  = \Io - \Jv %\\
         = \Io + J_{\rm{mv}} %\\
         = \Io + \Jv + 2 J_{\rm{mv}}
\end{aligned}
\end{equation*}
\end{remark}



\subsubsection{B: Flexure}

\begin{proposition}[Shear centre]
Any shear tangential stress field has a vanishing twisting moment about a pole of the cross section $\Omega$, called the shear centre $\mathrm{C}^{\rm sh}$ whose position is given by the formula:
$$
\begin{aligned}
\FrameBasis\times\bm{r}_{\mathrm{C}^{\rm sh}}
&=\BishearTensorRB^{-1} \int_{\Omega} G\mathbf{B}_{(b)}(\bm{r})^{\top}\, \FrameBasis\times \bm{r}\mathrm{~d}A \\
&=-\BishearTensorRB^{-1} \left(\int_{\Omega} G\mathbf{B}_{(b)}^{\top} [\bm{r}^{\times}]\mathrm{~d}A\right)\,\FrameBasis
\end{aligned}
$$
\end{proposition}

\cite{paradiso2019consistent}. 
Traditionally the evaluation of $\bm{\tau}_{\mathrm{sh}}$ is simplified by first evaluating the so called ``geometric" center of shear

\begin{equation}
\bm{r}_\psi=\FrameBasis \times \mathbf{J}_G^{-1} \bm{\gamma}_\psi=\FrameBasis \times \mathbf{J}_G^{-1} \int_{\Section} \mathbf{B}_{(b)} \FrameBasis?\times\bm{r} \mathrm{~d}A
\label{eq:para-47}
\end{equation}
See also \cite{elter1983two}.


\begin{Quote}
Regarding \eqref{eq:iesan-2-22} and the constant tensor relating to \(c_{\vartheta}\). 
The outward unit normal is $\bm{\nu}$ and the unit tangent is $\bm{t}$.
\[
\bm{J}
:=\int_{\Section}\bm{r}^\times\!\left(\nabla\bm{\WarpShearIesan})^{\!\top}
+\tfrac{\nu}{2}\,r^2\mathbf{1}\right)\,da
\;\in\mathbb{R}^{2}.
\]
% \paragraph{A boundary–area identity.}
Using the 2D identity \(\bm{r}\times\nabla\psi=\nabla\cdot(\psi\,\bm{r})\)
and the divergence theorem, one obtains, columnwise,
\begin{equation}
\bm{J}
= \oint_{\SectionBoundary}\!\Phi\,(\bm{r}\cdot\bm{t})\,ds
\;+\;\tfrac{\nu}{2}\int_{\Section}\! r^{2}\,\bm{r}^\perp\,da.
\label{eq:J-boundary-form}
\end{equation}

\subparagraph{Central symmetry about \(\CentroidLocation\).}
A planar set \(\Sigma\subset\mathbb R^2\) is \emph{centrally symmetric about \(\CentroidLocation\)} if it is invariant under the point reflection (half–turn)
\[
\mathcal R_{\CentroidLocation}:\ \bm{r}\ \mapsto\ 2\CentroidLocation-\bm{r}.
\]
Equivalently,
\[
\Sigma-\CentroidLocation \;=\; -\,(\Sigma-\CentroidLocation),
\qquad\text{ i.e. }\qquad
\bm s\in\Sigma-\CentroidLocation\ \Longrightarrow\ -\bm s\in\Sigma-\CentroidLocation,
\]
with the translated coordinates \(\bm s=\bm{r}-\CentroidLocation\).

If \(\partial\Sigma\) is piecewise \(C^1\), central symmetry about \(\CentroidLocation\) means the map
\[
\mathcal R_{\CentroidLocation}:\ \partial\Sigma\to\partial\Sigma,\qquad
\bm{r}\mapsto\bm{r}' = 2\CentroidLocation-\bm{r},
\]
is a bijection and, at paired points, the unit tangent and normal satisfy
\[
\bm t(\bm{r}')=-\,\bm t(\bm{r}),\qquad
\bm\nu(\bm{r}')=-\,\bm\nu(\bm{r}),\qquad
ds(\bm{r}')=ds(\bm{r}).
\]

\paragraph{Multiply connected sections.}
For a section with holes, \(\Sigma=\Sigma_{\text{outer}}\setminus\bigcup_{j} H_j\) is centrally symmetric about \(\CentroidLocation\) iff
\[
\mathcal R_{\CentroidLocation}(\Sigma_{\text{outer}})=\Sigma_{\text{outer}}
\quad\text{and}\quad
\{\mathcal R_{\CentroidLocation}(H_j)\}_j=\{H_j\}_j
\]
as sets; i.e., each hole is either individually centrally symmetric about \(\CentroidLocation\) or belongs to a centrally symmetric pair with respect to \(\CentroidLocation\).

\paragraph{Relation to mirrors.}
Two orthogonal reflection symmetries whose axes intersect at \(\CentroidLocation\) imply central symmetry about \(\CentroidLocation\) (but the converse need not hold: a shape can have central symmetry without any mirror symmetry).

\paragraph{Centroid remark.}
If \(\Section\) is centrally symmetric about \(\CentroidLocation\) and has uniform area density, then its area centroid coincides with \(\CentroidLocation\). 
The converse need not be true: \(\bar{\bm{r}}=\CentroidLocation\) does \emph{not} by itself imply central symmetry.

\paragraph{Vanishing Conditions}
The representation \eqref{eq:J-boundary-form} makes parity-based
cancellations immediate.
\begin{itemize}
\item \emph{Central symmetry about \(\CentroidLocation\)} and the zero-mean condition imply
\(\Phi\) is odd in \(\bm{r}\) while \(\bm{r}\cdot\bm{t}\) is even;
hence the boundary integral vanishes. The area integral
\(\int_{\Section} r^2\bm{r}^\perp\,da\) is odd and also vanishes.
Thus \(\bm{J}=\mathbf{0}\). This covers circles, ellipses, rectangles, regular
polygons, and multiply connected sections with centrally symmetric hole patterns.
\item \emph{Bisymmetry} (two orthogonal reflection axes through \(\CentroidLocation\))
kills both components of \(\bm{J}\) by the same parity argument.
\item If \(\nu=0\), the second term in \eqref{eq:J-boundary-form} disappears; a
single symmetry axis through \(\CentroidLocation\) is then sufficient to kill the
component perpendicular to that axis (bisymmetry kills both).
\end{itemize}


\paragraph{Summary.}
(i) The integral \(\bm{J}\) vanishes under the symmetry conditions listed above,
by \eqref{eq:J-boundary-form}. 
\end{Quote}


\[
\int_{\Section}\Big(\bm\WarpShearIesan\otimes(\bm{r}\times\mathbf i)-\bm{r}\otimes(\bm{r}\times\mathbf i)\Big)\mathrm{~d}A.
\]

\[
(\bm\WarpShearIesan-\bm{r})\otimes(\bm{r}\times\mathbf i)
=
\begin{bmatrix}
0\cdot 0 & 0 & 0\\
(\psi_y-y)\cdot 0 & (\psi_y-y)\,z & (\psi_y-y)\,(-y)\\
(\psi_z-z)\cdot 0 & (\psi_z-z)\,z & (\psi_z-z)\,(-y)
\end{bmatrix}.
\]

Therefore
\[
\int_{\Section}\Big(\bm\WarpShearIesan\otimes(\bm{r}\times\mathbf i)-\bm{r}\otimes(\bm{r}\times\mathbf i)\Big)\mathrm{~d}A
=
\displaystyle\int_{\Section} \begin{bmatrix}
0 & 0 & 0 \\
0 & z\,(\psi_y-y) & -y\,(\psi_y-y) \\
0 &  z\,(\psi_z-z)\mathrm{~d}A & -y\,(\psi_z-z)
\end{bmatrix}\mathrm{~d}A.
\]


\paragraph{psi - r}

\[
\int_{\Section}\Big(\bm\WarpShearIesan\otimes(\bm{r}\times\mathbf i)-\bm{r}\otimes(\bm{r}\times\mathbf i)\Big)\mathrm{~d}A.
\]

Equivalently, column-wise (since for any \( \bm a\otimes\bm b\), the (j)-th column is \(b_j \bm a\)):
\[
\int_{\Section} \bm\WarpShearIesan\otimes(\bm{r}\times\mathbf i) - \bm{r}\otimes(\bm{r}\times\mathbf i) \mathrm{~d}A
=
\int_{\Section} \Big[\,\bm 0\ \Big|\ z\,(\bm\WarpShearIesan-\bm{r})\mathrm{~d}A\ \Big|
\ -y\,(\bm\WarpShearIesan-\bm{r})\mathrm{~d}A,\Big],
\]

The nonzero \(yz\)-block is the \(2\times2\) matrix
\[
\mathbf I = \int_{\Section}\big(\bm{\WarpShearIesan}-\bm{r}\big)\otimes(\FrameBasis\times\bm{r})\mathrm{~d}A
=\underbrace{\int_{\Section}\bm{\WarpShearIesan}\otimes(\FrameBasis\times\bm{r})\mathrm{~d}A}_{\mathbf A}
-\underbrace{\int_{\Section}\bm{r}\otimes(\FrameBasis\times\bm{r})\mathrm{~d}A}_{\mathbf G = -\IntRoR\FrameBasis^\times}.
\]

Mixed term \(\mathbf A\) via integral theorems

Set \(f:=\tfrac12 r^2 \) so that \(\nabla f=\bm{r}\) and \(\FrameBasis\times\nabla f=\FrameBasis\times\bm{r}\).
Then
\[
\mathbf A=\int_{\Section}\bm{\WarpShearIesan}\otimes\FrameBasis\times\nabla f\mathrm{~d}A.
\]
Apply the rotated-gradient integration by parts (Stokes/Green in 2-D \eqref{eq:thm-stokes-green}) with \(\bm a=\bm\psi\), \(g=f\):
\[\mathbf A
=\oint_{\partial\Section} \tfrac12 r^2 \,\bm\psi\otimes(J^{\top}\bm\nu)\,\mathrm{d}s
-\int_{\Section} \tfrac12 r^2 \,(\nabla\bm\psi)\,J^{\top}\mathrm{~d}A.
\]

\[
\mathbf I
=\oint_{\partial\Section} \tfrac12 r^2 \,\bm\psi\otimes(J^{\top}\bm\nu)\,\mathrm{d}s
-\int_{\Section} \tfrac12 r^2 \,(\nabla\bm\psi)\,J^{\top}\mathrm{~d}A
-\mathbf M\,J^{\top}.
\]


\paragraph{Warping Integrals}.
\begin{Quote}

Setting \(\bm w=\bm\WarpShearIesan\) with the BVP data in Green's identity \eqref{eq:thm-green-vec} with the Divergence Theorem \eqref{eq:div-thm-vec}  gives
\[
\int_{\Section}\nabla\bm\WarpShearIesan(\nabla\bm\WarpShearIesan)^{\!\top}\mathrm{~d}A
=2\int_{\Section}\bm\WarpShearIesan\otimes(\bm{r}-\CentroidLocation)\mathrm{~d}A
+\int_{\SectionBoundary}\bm\WarpShearIesan\otimes
\Big(\big(\operatorname{dev}\bm{r}\otimes\bm{r}-\bm{r}\otimes\CentroidLocation\big)^{\!\top}\bm\nu\Big)\,\mathrm{d}s,
\]
while
\[
-\int_{\Section}\nabla\bm\WarpShearIesan\mathrm{~d}A
=-\int_{\SectionBoundary}\bm\nu\otimes\bm\WarpShearIesan\,\mathrm{d}s.
\]

Therefore, the identity
\[
\int_{\Section}\nabla\bm\WarpShearIesan(\nabla\bm\WarpShearIesan)^{\!\top}\mathrm{~d}A
= -\int_{\Section}\nabla\bm\WarpShearIesan\mathrm{~d}A
\]
holds \emph{if and only if}
\[
\int_{\SectionBoundary}\!\!\Big[
\bm\WarpShearIesan\otimes\Big(\big(\operatorname{dev}\bm{r}\otimes\bm{r}-\bm{r}\otimes\CentroidLocation\big)^{\!\top}\bm\nu\Big)
+\bm\nu\otimes\bm\WarpShearIesan\Big]\mathrm{d}s
\;+\;2\int_{\Section}\bm\WarpShearIesan\otimes(\bm{r}-\CentroidLocation)\mathrm{~d}A
\;\equiv\;\mathbf{0}.
\tag{\(*\)}
\]
This is the necessary and sufficient condition.

From \((*)\) you can read off a few \emph{simple sufficient scenarios}:
\begin{itemize}
  \item If \(\displaystyle \int_{\Section}\bm\WarpShearIesan\otimes(\bm{r}-\CentroidLocation)\mathrm{~d}A=\mathbf{0}\) (i.e., \(\bm\WarpShearIesan\) is \(L^2\)-orthogonal to \(\bm{r}-\CentroidLocation\)) \emph{and}
  \[
  \big(\operatorname{dev}\bm{r}\otimes\bm{r}-\bm{r}\otimes\CentroidLocation\big)^{\!\top}\bm\nu=-\,\bm\nu
  \quad\text{on }\SectionBoundary,
  \]
  then the boundary integral in \((*)\) vanishes and the equality holds.
  \item On a centrally symmetric \(\Section\) about \(\CentroidLocation\), if \(\bm\WarpShearIesan\) is odd with respect to \(\CentroidLocation\) (so the area integral in \((*)\) vanishes) \emph{and} the boundary factor
  \(\big(\operatorname{dev}\bm{r}\otimes\bm{r}-\bm{r}\otimes\CentroidLocation\big)^{\!\top}\bm\nu+\bm\nu\)
  is pointwise odd along symmetric boundary pairs, then the boundary integral also vanishes.
\end{itemize}

A final algebraic remark: since \(\int_{\Section}\nabla\bm\WarpShearIesan(\nabla\bm\WarpShearIesan)^{\!\top}\) is symmetric positive-semidefinite, the matrix \(-\int_{\Section}\nabla\bm\WarpShearIesan\) must itself be symmetric positive-semidefinite whenever the identity holds. Equivalently,
\[
\int_{\Section}\nabla\bm\WarpShearIesan\ \text{ must be symmetric and negative-semidefinite.}
\]
This is a necessary condition that you can check independently via \(\int_{\Section}\nabla\bm\WarpShearIesan=\int_{\SectionBoundary}\bm\nu\otimes\bm\WarpShearIesan\,\mathrm{d}s\).

\end{Quote}

